\documentclass[10pt,a4paper]{amsart}
\usepackage[margin=19mm]{geometry}
\usepackage{amsmath,amssymb,mathtools}
\usepackage[scr=rsfs,cal=euler]{mathalfa}
\usepackage{xfrac}
\usepackage[unicode,hidelinks,bookmarks=false]{hyperref}
\newcommand{\ie}{i.e.\ }
\newcommand{\cf}{cf.\ }
\newcommand{\resp}{respectively\ }
\newcommand{\Subsection}{\subsection}
\providecommand{\nobreakdash}{\nobreak\hbox{-}}
\setlength{\emergencystretch}{2em}
\multlinegap0pt
\usepackage{amssymb}
\usepackage{dsfont}
\usepackage[shortlabels]{enumitem}
\usepackage{xcolor}
\newtheorem{proposition}{Proposition}
\newcommand{\map}[1]{\textbf{#1}}
\title{Asymptotic normality of pattern counts in random maps II}
\author{Eva-Maria Hainzl}
\date{Source: arXiv:2603.19485v1 (CC BY 4.0)}
\begin{document}
\maketitle

\noindent\emph{Source and licence.} Asymptotic normality of pattern counts in random maps II. Version arXiv:2603.19485v1 (CC BY 4.0), distributed by arXiv under the Creative Commons Attribution 4.0 International licence, http://creativecommons.org/licenses/by/4.0/, as recorded in the arXiv licence metadata for this exact version and shown on the abstract page https://arxiv.org/abs/2603.19485v1. Full licence notice: \texttt{licenses/arxiv-2603.19485-CC-BY-4.0.txt}.\par\smallskip

\noindent\textit{Editorial scope.} Counted unit: the article's proposition thm:nonsep\_dec (source0.tex lines 438--446). The article's complete original proof of it is delivered with it (source0.tex lines 448--474, one \texttt{proof} environment of 27 lines). No mathematics is written, repaired or re-derived: the statement and the proof are the article's own lines, in the article's own notation, and no retained line is substituted or re-flowed.  No further source-local context is needed: every cross-reference of every retained line resolves inside the two delivered spans.  One declared adaptation: exactly 1 ancillary strings inside the retained spans are omitted (image-inclusion commands and, where the source repeats a label, the repeated label definitions) and no other line differs from the source; the figure environments, with their placement, captions and labels, are kept, and the package records the omitted strings in fidelity receipts bound to the affected bodies.  No retained line contains a citation, so the wrapper carries no bibliography and no bibliography entry.  The retained lines contain the article's own diagram code, which the wrapper typesets with the standard tikz packages; no image file is delivered and the package declares no asset dependency.  Editorial formatting only, in addition: an AMS \texttt{amsart} preamble that restates the article's own declarations and macros used by the retained lines, namely the environments \texttt{proposition} on separate counters, together with the standard packages those lines need.  The retained environments print in this document as Proposition 1 in order of appearance, whereas the article prints the same items with the numbering of its own full document.  The complete native source of the article as distributed in the arXiv e-print for this exact version is bundled unchanged as \texttt{sources/196776/source0.tex}.
\begin{proposition} \label{thm:nonsep_dec}
    Let $\map{p}$ be a planar 2-connected map with $e$ edges, a boundary of length $v$ and $r$ rotational symmetries. Further, we enumerate its intersection types by $1$ to $I$, and let $e_i$ be the number of edges, $v_i$ the valency of the root face, $d_{i,j}$ the number of deep faces with valency $j$ and $r_i$ the number of rotations of intersection type $i \in [I]$. Then the generating function $N(z,u,x)$ for planar 2-connected maps with some distinguished labeled pattern occurrences of $\map{p}$, intersecting at most pairwise satisfies the equation
    \begin{align*}
        N(z,u,x) &= zu\frac{zu+\frac{uN(z,1,x)-N(z,u,x)}{1-u}}{1-\left(zu+\frac{uN(z,1,x)-N(z,u,x)}{1-u}\right)} \\
        &\quad + x\frac{rz^{e-v+1}}{u^{v-2}}F_{v-2}(z,u,N(z,u,x),n_2(z,x),\dots,n_{v-2}(z,x))\\
        &\quad + x^2\sum_{i=0}^I \frac{r_iz^{e_i-v_i+1}}{u^{v_i-2}}F_{v_i-2}(z,u,N(z,u,x),n_2(z,x),\dots,n_{v_i-2}(z,x))\prod_{j\geq 2} n_j(z,x)^{d_{i,j}}
    \end{align*}
    where all $F_v(z,u,y_0,y_1,\dots, y_{v-1}), v>0$ are polynomials.
\end{proposition}

\begin{proof}
    Analogous to the bivariate case, the first term in the equation counts all maps where the root edge is not incident to a pattern occurrence and the rest corresponds to the cases where the root edge is incident to a labeled pattern occurrence or an intersection type that we isolate.
    Next, let us point out that in case that the root edge is incident to a labeled pattern occurrence, the map without the pattern occurrence (or the occurrence of the pair of intersecting patterns) consists of a non-empty sequence of 2-connected maps or single edges as illustrated in Figure~\ref{fig:non-sep}. 

    \begin{figure}
        \centering
        
        \caption{Decomposition of 2-connected maps}
        \label{fig:non-sep}
    \end{figure}
    
    To describe those maps, we introduce a variable $w$ that marks the number of edges incident to the pattern occurrence in addition to $z$ marking the number of edges and $u$ marking the root face valency. Each of the 2-connected maps in this sequence contribute at least one edge to the boundary of the pattern occurrence(s) and at most one less than their root face valency. Therefore, the generating function of each of those maps equals
    \begin{align*}
        \sum_{m\geq 2} n_m(z,x)u^m(w+w^2+\dots+w^{m-1})
        &=\sum_{m\geq 2} n_m(z,x)u^m\frac{w-w^m}{1-w}\\
        &=\frac{wN(z,u,x)-N(z,uw,x)}{1-w}
    \end{align*}
    Thus, the generating function of a non-empty sequence of edges or 2-connected maps that contribute $b-2$ edges to the boundary of a pattern occurrence or an intersection type is described by
    \begin{align*}
        [w^{b-2}] \frac{zu^2w+\frac{wN(z,u,x)-N(z,uw,x)}{1-w}}{1-\left(zu^2w+\frac{wN(z,u,x)-N(z,uw,x)}{1-w}\right)}
    \end{align*}
    which obviously is a polynomial $F_{b-2}(z,u,y_0,y_1,\dots, y_{b-3})$ in $z,u,$ $y_0=N(z,u,x)$ and $y_1 =n_2(z,x), \dots, y_3 = n_{b-2}(z,x)$.

    Subsequently, we add $e-v+1$ edges of the pattern occurrence which decreases the root face valency by $v-2$. Further, we multiply the expression by the number of rotational symmetries of the pattern to account for all possible ways of attaching the pattern occurrence to the sequence of 2-connected maps and mark the term with $x$.

    In the case where the root edge is incident to an intersection type, the above argument holds as well. Additionally, the deep faces are filled by 2-connected maps with suitable root face valency (which have a simple boundary by definition).
\end{proof}

\end{document}
